% =====================================================================================
% SEVA_v8.tex --- reconciled manuscript draft (v8). IEEE conference style (drop-in for v7.1.3).
% Built 2026-06-01 from the RECONCILE-v713 finalized claim set (PAPER_EDITS_LOG.md, commit 8becc5d).
% -------------------------------------------------------------------------------------
% SOURCE TAGS: every quantitative / empirical claim carries a "% [TAG]" margin marker keyed
%   to a result entry in PAPER_EDITS_LOG.md. These are author-verifiable provenance markers;
%   STRIP ALL "% [...]" markers before submission (grep '% \[').
%
% TAG LEGEND (result -> claim it licenses):
%   [E-CAL-1]          frozen, non-oracle, held-out calibration: templated 0% ASR @ ~0.58% in-domain DocFPR
%   [SEEDS-1]          3-seed stability (0% across seeds 42/7/123; complete per-seed FPR)
%   [STEP3]            in-domain baseline gate: cluster_coh gap/SNR, L1 0% ASR, DocFPR 0.40-0.52%, CUDA latency 13.7-15.7 ms
%   [E2-2]             in-domain soft-signal SNR collapse (kw_density 38->8) -> composite-as-ablation evidence
%   [E-CAL-2]          composite under keyword-dropping adaptive adversary collapses to 49-57% ASR
%   [CHEAP-MUST-1]     cluster_coh hard gate 100/98/100% (templated/PoisonedRAG/L2-L3) @ 0.69% DocFPR; >=2-aggregation query-FPR 3.15%->0.90%
%   [PR-GATE-1]        real RELEASED PoisonedRAG on Security-SE: cluster_coh 98%, composite 72%
%   [PR-XDOMAIN]       RELEASED PoisonedRAG on Natural Questions: cluster_coh 82%, MinHash 0%, s_nd 52% @ matched FPR
%   [PR-XDOMAIN-HOTPOT] RELEASED PoisonedRAG on HotpotQA: cluster_coh 97%, MinHash 0%, s_nd 98%
%   [E4-HH]            reproduced matched-FPR head-to-head vs RAGDefender: SEVA 100% vs ~89%; RAGDefender native ~51% DocFPR
%   [OPEN-CAL-1]       non-oracle calibration protocol (calibrate on a reference attack sample, deploy frozen)
%   [LAT-M4]           Apple M4 latency files: ~32-42 ms mean (supersedes v7.1.3's 26 ms; sub-2 ms is dead)
%
% SCOPE DISCIPLINE (held exactly): attack scope = "templated / multi-passage / clustered injection
%   (PoisonedRAG family)" --- NEVER "near-duplicate corpus poisoning" as an umbrella claim. Dedup
%   claims scoped to LEXICAL filtering; s_nd reported as a competitive embedding-dedup baseline
%   (cluster_coh >= s_nd: edge on NQ 82% vs 52%, parity elsewhere) --- never "beats all dedup".
%   Latency only the measured ~13-16 ms GPU / ~32-42 ms M4 numbers.
%
% CHECKPOINT STATUS:
%   [Ckpt 1 -- THIS FILE] frontmatter + abstract + contributions C1-C6 + one-line section outline. No prose bodies.
%   [Ckpt 2] Intro + Related Work + Threat Model bodies (Sec I-III).
%   [Ckpt 3] Architecture + Evaluation bodies + reworked tables (Sec IV-V).
%   [Ckpt 4] Discussion + Limitations + Conclusion bodies (Sec VI-VIII).
% =====================================================================================

\documentclass[conference]{IEEEtran}
\IEEEoverridecommandlockouts
\usepackage{cite}
\usepackage{amsmath,amssymb,amsfonts}
\usepackage{algorithm}
\usepackage{algorithmic}
\usepackage{graphicx}
\usepackage{textcomp}
\usepackage{xcolor}
\usepackage{booktabs}
\usepackage{multirow}
\usepackage{url}

\begin{document}

\title{SEVA: Lightweight, LLM-Free Detection of Templated\\
Corpus Poisoning in Retrieval-Augmented Generation}

\author{\IEEEauthorblockN{Anonymous Author(s)}
\IEEEauthorblockA{Manuscript submitted for blind review}}

\maketitle

% =====================================================================================
% ABSTRACT --- reconciled set (RECONCILE-v713). Deployability-first, cluster_coh-centric.
% Every quantitative clause carries its source tag at line end.
% =====================================================================================
\begin{abstract}
Retrieval-Augmented Generation (RAG) systems are vulnerable to corpus
poisoning---adversarial document injection that corrupts retrieval without
modifying model weights or triggering lexical filters. Existing defenses
require LLM decoder access, multiplicative inference overhead, or cryptographic
pre-registration, making them impractical for the on-device, offline deployment
that local AI increasingly demands. We present SEVA, a fully local, LLM-free
detector built on a single geometric signal---per-document $K$-nearest-neighbour
cluster coherence (\texttt{cluster\_coh})---operated as a hard gate. SEVA detects
templated multi-passage poisoning, the dominant published attack pattern
(PoisonedRAG), at 0\% attack-success and 0.58\% document-level false-positive
rate % [E-CAL-1, SEEDS-1]  (confirm exact DocFPR vs STEP3 0.45-0.52% at table-build, Ckpt 3)
under density-agnostic, non-oracle (frozen, held-out) calibration, validated
across three seeds and 1--10\% contamination density. % [E-CAL-1, SEEDS-1, STEP3]
On the attack authors' own released PoisonedRAG poison, SEVA catches the attack
cross-domain at 82--98\% (Security Stack Exchange, Natural Questions, and
HotpotQA) under non-oracle calibration. % [PR-GATE-1, PR-XDOMAIN, PR-XDOMAIN-HOTPOT]
We further show that lexical duplicate filtering---the method PoisonedRAG itself
evaluated and found insufficient---is corpus-fragile: at matched false-positive
rate it fails entirely on duplicate-dense benchmarks where cluster coherence
holds. % [PR-XDOMAIN, PR-XDOMAIN-HOTPOT]
Against a keyword-dropping adaptive adversary, a multi-signal composite collapses
to 49--57\% attack-success while the geometric hard gate holds 0\%; % [E-CAL-2, CHEAP-MUST-1]
distilling to the geometric core is what makes SEVA both robust and deployable.
SEVA runs in approximately 13--16\,ms on a commodity GPU and 32--42\,ms on an
Apple-Silicon laptop, % [STEP3, LAT-M4]
with no LLM, API, or pre-registration; in a reproduced, shared-corpus,
matched-FPR head-to-head it catches 100\% of templated poison versus the
per-query state of the art's 89\%. % [E4-HH]
SEVA targets structured multi-passage injection; single-document injections
crafted to mimic the clean distribution are out of scope.
\end{abstract}

\begin{IEEEkeywords}
Retrieval-Augmented Generation, Corpus Poisoning, Anomaly Detection, Embedding
Security, Geometric Detection, On-Device Security.
\end{IEEEkeywords}

% =====================================================================================
\section{Introduction}

\subsection{Motivation}
Retrieval-Augmented Generation (RAG) grounds a language model's output in an
external document corpus, decoupling parametric model knowledge from a
non-parametric, operator-maintained vector store. That decoupling is also RAG's
security liability: it relocates factual authority to a corpus an adversary may
be able to write to. \emph{Corpus poisoning} exploits this. An adversary with
write access to the ingestion pipeline injects crafted documents that are
retrieved at query time and steer the model toward attacker-chosen outputs---%
operating at the infrastructure layer beneath the LLM, without modifying model
weights, and without leaving the lexical traces a keyword filter would catch.

The dominant instantiation is PoisonedRAG~\cite{poisonedrag}, which decomposes
each malicious text into a retrieval-satisfying component and a
generation-satisfying payload and injects roughly five such passages per target
query, reaching 90--97\% attack-success even in corpora of millions of
documents. Two properties of this attack are decisive for the defender. First,
it is \emph{templated}: the passages injected for a target query are generated
from a shared recipe and are therefore mutually similar. Second, it is
\emph{multi-passage}: several co-retrieved passages are injected per query to
dominate the retrieval window. Together they mean the poison enters the corpus
as a tight \emph{cluster} of mutually similar documents---a geometric signature
intrinsic to the attack's own mechanism, and, as we show, present regardless of
how much of the corpus is poisoned. % [STEP3]

\subsection{The Deployability Gap and Our Approach}
The defenses that drive standard attack-success below roughly 20\% each carry a
deployment cost that an on-device, offline RAG system cannot pay. Certified
defenses such as RobustRAG~\cite{robustrag} incur multiplicative---reportedly
5$\times$---LLM inference overhead. The Attention-Variance Filter~\cite{avfilter}
requires decoder attention access and still faces 35\% attack-success under an
adaptive adversary. Cryptographic provenance schemes such as
RAGShield~\cite{ragshield} require re-registration of the entire corpus. The
closest peer to our setting, RAGDefender~\cite{ragdefender}, is an LLM-free
post-retrieval filter, and is the baseline we reproduce head-to-head. As local
and laptop-class RAG proliferates, a defense must increasingly run where the
corpus lives---on commodity or Apple-Silicon hardware, offline, with no LLM in
the loop and no per-query model call.

SEVA meets exactly that envelope. It detects the templated multi-passage attack
with a single per-document geometric signal---the local $K$-nearest-neighbour
cluster coherence of the corpus---operated as a hard gate, with no LLM call, no
external API, and no corpus pre-registration. On commodity GPU hardware it runs
in approximately 13--16\,ms per query, and on an Apple-Silicon laptop in
32--42\,ms. % [STEP3, LAT-M4]
The geometric core is what survives an adaptive adversary: a multi-signal
composite that also reads soft linguistic statistics is domain-confounded and
collapses under a keyword-dropping attacker, whereas the cluster-coherence gate
is immune to soft-signal evasion by construction. % [E2-2, E-CAL-2, CHEAP-MUST-1]

We evaluate SEVA under a deliberately demanding protocol: density-agnostic,
non-oracle (frozen, held-out) calibration; an in-domain corpus that removes the
domain shortcut most defense evaluations enjoy; contamination densities of
1--10\% with complete per-seed false-positive reporting; the attack authors' own
released PoisonedRAG poison across three corpora; and a reproduced,
shared-corpus, matched-FPR head-to-head against the per-query state of the
art. % [E-CAL-1, SEEDS-1, PR-XDOMAIN, E4-HH]
Across this protocol SEVA holds 0\% attack-success on templated poison in-domain
and catches the released attack cross-domain at 82--98\%, while lexical
duplicate filtering---the defense PoisonedRAG itself tested---fails entirely at
matched false-positive rate on duplicate-dense corpora. % [E-CAL-1, PR-XDOMAIN, PR-XDOMAIN-HOTPOT]

\subsection{Contributions}
SEVA makes the following contributions. Each empirical claim is traceable to a
released result and scoped exactly to what was measured.

\begin{itemize}
\item \textbf{C1 --- Deployability.} A fully local, LLM-free corpus-poisoning
detector that runs in approximately 13--16\,ms on a commodity GPU and 32--42\,ms
on an Apple-Silicon laptop, % [STEP3, LAT-M4]
with no LLM call, external API, or cryptographic pre-registration---a deployment
profile that LLM-based, decoder-access, and cryptographic defenses cannot match.

\item \textbf{C2 --- Cross-domain detection of the dominant attack.} SEVA detects
templated multi-passage poisoning at 0\% attack-success and 0.58\% document-level
false-positive rate under frozen, non-oracle, three-seed calibration, % [E-CAL-1, SEEDS-1]
and catches the attack authors' own released PoisonedRAG poison cross-domain at
82--98\% across three corpora (Security Stack Exchange, Natural Questions, and
HotpotQA). % [PR-GATE-1, PR-XDOMAIN, PR-XDOMAIN-HOTPOT]

\item \textbf{C3 --- What cluster coherence adds over embedding deduplication.}
The detection signal is a single geometric gate that (i)~is density-invariant
across 1--10\% contamination---a property pairwise deduplication lacks; % [STEP3]
(ii)~matches or beats embedding near-duplicate detection (an edge on Natural
Questions, 82\% vs.\ 52\%; parity elsewhere) by keying on the multi-passage
cluster rather than a single nearest neighbour; % [PR-XDOMAIN, PR-XDOMAIN-HOTPOT]
and (iii)~catches the attack where lexical duplicate filtering fails at matched
false-positive rate. % [PR-XDOMAIN]

\item \textbf{C4 --- Adaptive robustness through distillation.} We show that a
multi-signal, SNR-weighted composite is domain-confounded and collapses to
49--57\% attack-success under a keyword-dropping adversary, while the geometric
core operated as a hard gate is invariant to soft-signal evasion (0\%); % [E-CAL-2, CHEAP-MUST-1]
the composite is reported as the ablation that motivates distilling SEVA to its
geometric core.

\item \textbf{C5 --- Methodological rigor.} SEVA is evaluated under
density-agnostic, non-oracle (frozen, held-out) calibration % [OPEN-CAL-1, E-CAL-1]
at 1--10\% contamination density with complete per-seed false-positive
reporting, % [SEEDS-1]
and against a reproduced---not number-lifted---matched-FPR head-to-head with the
LLM-free per-query state of the art (SEVA 100\% vs.\ 89\% on templated poison; the
baseline's native operating point removes roughly half of clean documents
in-domain). % [E4-HH]

\item \textbf{C6 --- Lexical deduplication is corpus-fragile.} We report a new
finding: lexical duplicate filtering---the defense PoisonedRAG itself tested---is
effective only on lexically sparse corpora and fails at matched false-positive
rate on duplicate-dense standard benchmarks, where semantic cohesion continues to
detect the attack. % [PR-XDOMAIN, PR-XDOMAIN-HOTPOT]
\end{itemize}

% =====================================================================================
\section{Related Work}

\subsection{Corpus Poisoning Attacks}
PoisonedRAG~\cite{poisonedrag} established the dominant formalism for
knowledge-corruption attacks, decomposing each malicious text into a
retrieval-satisfying and a generation-satisfying component and reaching 90--97\%
attack-success with five injected passages in corpora of up to millions of
documents. CatPoison~\cite{catpoison} lifts instance-level attacks to the
category level by anchoring poisoned embeddings near category-query centroids
(mean 80.1\% attack-success). AutoTPA~\cite{autotpa} automates trigger-targeted
poisoning, reporting 95.5\% retrieval success. SEVA targets the templated,
multi-passage family these attacks define---the regime in which poison enters
the corpus as a cluster of mutually similar passages.

\subsection{Defenses and Their Deployment Cost}
Detectors that do not call an LLM have been limited in power: perplexity-based
detection reaches only AUC 0.25--0.30 against PoisonedRAG, because adversarially
optimized text can be kept within the normal perplexity
range~\cite{poisonedrag}, and paraphrasing defenses degrade legitimate
retrieval. The defenses that do drive standard attack-success low pay for it at
deployment. RobustRAG~\cite{robustrag} provides certified robustness but incurs
roughly 5$\times$ LLM inference overhead. The Attention-Variance
Filter~\cite{avfilter} requires decoder attention access and reaches 35\%
attack-success under an adaptive GCG-Poison adversary. RAGShield~\cite{ragshield}
achieves 0\% attack-success against external injection but requires cryptographic
re-registration of the corpus and still exhibits 17.5\% insider attack-success.
RAGDefender~\cite{ragdefender} is the closest peer to SEVA---an LLM-free
post-retrieval ML filter reporting roughly 7.5--12.5\% attack-success at about
1\% density, and, like SEVA, it evaluates an adaptive adversary of its own. We
take RAGDefender as our head-to-head baseline and reproduce it on a shared corpus
and attack at matched false-positive rate. % [E4-HH]
SEVA occupies the gap none of these fill: a corpus-level, LLM-free detector with
no decoder access, no pre-registration, and millisecond latency on laptop-class
hardware.

\subsection{Lexical and Semantic Duplicate Filtering}
PoisonedRAG itself evaluated duplicate-text filtering as a candidate defense and
reported it insufficient~\cite{poisonedrag}. Lexical near-duplicate
detectors---MinHash- and SimHash-style filters that flag verbatim or
near-verbatim overlap% TODO Ckpt3: \cite{minhash,simhash}
---are attractive because they are cheap and LLM-free, but their effectiveness is
a property of the corpus, not of the attack: a corpus already rich in benign
duplication forces the lexical threshold so high that injected passages slip
beneath it. We make this precise in Section~\ref{sec:eval}: at matched
false-positive rate, lexical filtering catches the templated attack on lexically
sparse corpora and fails entirely on duplicate-dense standard benchmarks, where
the cluster-coherence signal continues to detect it. % [PR-XDOMAIN, PR-XDOMAIN-HOTPOT]
We additionally report embedding-based near-duplicate detection (maximum cosine
to the nearest corpus neighbour) as a competitive baseline; cluster coherence
matches or exceeds it, with a clear edge on Natural Questions. % [PR-XDOMAIN]

% =====================================================================================
\section{Threat Model}
\label{sec:threat}

\subsection{System and Adversary Model}
Let $D=\{d_1,\dots,d_N\}$ be the corpus and $f:\mathcal{T}\to\mathbb{R}^{d}$ the
embedding function; a RAG system retrieves the top documents for a query $q$ and
conditions generation on them. The adversary has write access to the ingestion
pipeline and injects a set of crafted passages $D_{\mathrm{adv}}\subset D$ such
that, for a set of target queries $Q_{\mathrm{adv}}$, the system retrieves
elements of $D_{\mathrm{adv}}$ with high probability and is steered to an
attacker-specified output. The adversary is black-box with respect to the
defense: it knows the (public) embedding architecture but not SEVA's thresholds,
calibration data, or operating point. Following the released PoisonedRAG
construction~\cite{poisonedrag}, each target query receives $V=5$ injected
passages, each combining a retrieval-satisfying component with a
generation-satisfying payload.

\subsection{Attack Scope}
SEVA targets \emph{templated, multi-passage, clustered injection}---the
PoisonedRAG family---in which an adversary injects multiple co-retrieved
passages, generated from a shared recipe, to dominate the retrieval window for a
target query. This is the dominant corpus-poisoning pattern in the literature,
and its defining property is geometric: the injected passages form a tight
cluster of mutually similar embeddings.

\subsection{Adaptive Adversary}
We additionally consider an adaptive adversary that knows soft statistical
detectors exist and normalizes them away---suppressing the elevated keyword
density and regularizing the sentence-level structure that a linguistic signal
would key on. This is the adversary class that defeats keyword- and
structure-based detection, and that RAGDefender~\cite{ragdefender} likewise
evaluates. In our evaluation it corresponds to the conditions that ablate the
soft signals (Section~\ref{sec:eval}); SEVA's geometric core is defined so that
this adversary has nothing to normalize. % [E-CAL-2]

\subsection{Defense Objective}
Writing $\mathrm{ASR}=\mathrm{FN}/(\mathrm{TP}+\mathrm{FN})$ and
$\mathrm{FPR}=\mathrm{FP}/(\mathrm{FP}+\mathrm{TN})$, SEVA minimizes
attack-success subject to a false-positive ceiling
$\mathrm{FPR}_{\mathrm{TARGET}}=0.69\%$, using only benign queries for
calibration---no LLM access, no labeled poison examples, and no oracle knowledge
of the poison density. % [OPEN-CAL-1]

\subsection{Out of Scope}
A single document crafted to statistically mimic the clean corpus distribution
exhibits no clustered cohesion and falls outside SEVA's geometric regime; we
treat it as a stated residual of corpus-level detection
(Section~\ref{sec:limits}) rather than a target of the present work.

% =====================================================================================
\section{SEVA Architecture}
% -------------------------------------------------------------------------------------
% OUTLINE (bodies + Algorithm at Ckpt 3):
%   A. The geometric core -- cluster_coh = per-doc mean pairwise cosine over K=5 nearest CORPUS neighbours;
%      precomputed corpus-level statistic, immune to per-query retrieval contamination; K_FETCH=20 over-fetch->rerank. [STEP3]
%   B. The hard gate -- flag a document iff cluster_coh exceeds a calibrated tau; templated/clustered poison sits
%      ~0.99 vs clean ~0.73 (gap ~0.235-0.247 in-domain, SNR ~5.8-6.0). [STEP3]
%   C. >=2-per-query aggregation -- flag a QUERY only when >=2 of its retrieved docs gate -> query-FPR 3.15%->0.90%
%      at zero catch cost. [CHEAP-MUST-1]
%   D. Density-agnostic non-oracle calibration -- tau set at the (1-FPR_TARGET) clean-score quantile via 50-iter
%      binary search on a reference attack sample, deployed FROZEN; no labeled poison, no oracle density. [OPEN-CAL-1, E-CAL-1]
%   E. The composite as ablation -- the 10-signal SNR-weighted fusion is presented as the ABLATION that motivates
%      distillation: in-domain the soft linguistic signals are domain-confounded (kw_density SNR 38->8) and collapse
%      under keyword-dropping; the geometric core, hard-gated, holds. Frame as honest evolution-through-rigor. [E2-2, E-CAL-2]
%   F. LLM-free CPU pipeline (core-identity invariant).
% -------------------------------------------------------------------------------------

% =====================================================================================
\section{Evaluation}
\label{sec:eval}
% -------------------------------------------------------------------------------------
% OUTLINE (bodies + reworked tables at Ckpt 3):
%   A. Setup -- in-domain Security-SE 100k deduped (deliberate domain-confound control); RELEASED PoisonedRAG
%      poison on NQ + HotpotQA; bge-large-en-v1.5 (1024-d, normalized); FAISS HNSW (M=32, efC=200); seeds 42/7/123;
%      FPR_TARGET 0.69%. [STEP3, PR-XDOMAIN]
%   B. In-domain frozen detection (HEADLINE) -- templated 0% ASR @ ~0.58% DocFPR, frozen non-oracle held-out,
%      3-seed, across 1-10% density. [E-CAL-1, SEEDS-1, STEP3]  -> REWORKED Table (was Table V; in-domain)
%   C. Cross-domain catch on released PoisonedRAG -- Security-SE 98% / NQ 82% / HotpotQA 97% @ 0.69% DocFPR,
%      non-oracle. [PR-GATE-1, PR-XDOMAIN, PR-XDOMAIN-HOTPOT]  -> NEW 3-corpus cross-domain Table
%   D. Lexical dedup is corpus-fragile -- MinHash 0% at matched FPR on duplicate-dense NQ/HotpotQA where
%      cluster_coh holds 82-97%; embedding-dedup s_nd competitive (52% NQ / 98% HotpotQA) => cluster_coh >= s_nd
%      (edge NQ, parity elsewhere). [PR-XDOMAIN, PR-XDOMAIN-HOTPOT]
%   E. Matched-FPR head-to-head vs RAGDefender -- SEVA 100% vs ~89% on templated; RAGDefender native ~51% DocFPR
%      (no no-attack gate -> undeployable in-domain). [E4-HH]  -> REWORKED Table (was Table IX; real, not number-lifted)
%   F. Hard gate + aggregation + adaptive -- cluster_coh hard gate 100/98/100% (templated/PoisonedRAG/L2-L3)
%      @ 0.69% DocFPR; composite collapses L2/L3 to 49-57% while the gate holds 0%. [CHEAP-MUST-1, E-CAL-2]
%   G. Efficiency -- ~13-16 ms mean CUDA (RTX 5080, p95 <=18.8 ms) / ~32-42 ms Apple M4; CPU; FAISS HNSW O(log N).
%      [STEP3, LAT-M4]  -> REWORKED latency Table (was Table IV; corrected M4, sub-2 ms killed)
% -------------------------------------------------------------------------------------

% =====================================================================================
\section{Discussion}
% -------------------------------------------------------------------------------------
% OUTLINE (bodies at Ckpt 4):
%   A. Deployment implications -- laptop / Apple-Silicon, offline, no LLM/API/crypto; corpus-level signal the
%      per-query SOTA structurally cannot replicate. [LAT-M4, STEP3]
%   B. Distillation through rigor -- the gauntlet (adaptive + in-domain control) revealed the soft signals were
%      domain-confounded and adaptively fragile; we distilled to the geometric core. The composite is the ablation. [E2-2, E-CAL-2]
%   C. Cluster coherence vs embedding deduplication -- density-invariance + cluster structure -> edge on NQ; report
%      s_nd as a competitive baseline, not a defeated one. [PR-XDOMAIN]
% -------------------------------------------------------------------------------------

% =====================================================================================
\section{Limitations}
\label{sec:limits}
% -------------------------------------------------------------------------------------
% OUTLINE (bodies at Ckpt 4 -- scoped boundaries stated with confidence, NOT a wound):
%   1. Single-document diffuse injection -- an adversary injecting a single document crafted to mimic the clean
%      distribution exhibits no clustered cohesion and is outside SEVA's geometric scope (one boundary sentence;
%      RobustRAG/AV-Filter-style residual). [private record stays private]
%   2. Calibration assumption -- like all calibrated detectors, SEVA assumes a reference sample of the attack
%      class; we show frozen reference-calibration generalizes to held-out poison. [E-CAL-1]
%   3. Two-point corpus scaling -- validated at 10k and 100k only. [Lim5]
%   4. Single encoder -- all results use bge-large-en-v1.5; encoder generalization is future work. [Lim7]
% -------------------------------------------------------------------------------------

% =====================================================================================
\section{Conclusion}
% -------------------------------------------------------------------------------------
% OUTLINE (body at Ckpt 4):
%   Restate the scoped thesis: an LLM-free geometric detector of templated multi-passage poisoning, deployable at
%   ~13-42 ms across commodity GPU and Apple-Silicon laptop hardware, catching the released PoisonedRAG attack
%   cross-domain at 82-98%, with lexical duplicate filtering shown corpus-fragile and a reproduced matched-FPR
%   head-to-head against the per-query SOTA. Future work: encoder generalization; density beyond 10%.
%   [E-CAL-1, PR-XDOMAIN, E4-HH]
% -------------------------------------------------------------------------------------

% =====================================================================================
% REFERENCES --- carried over from v7.1.3 (drop-in). TODO Ckpt 3: add MinHash (Broder 1997) +
%   SimHash (Charikar 2002) for the lexical-dedup-fragility claim; add Natural Questions
%   (Kwiatkowski et al. 2019) + HotpotQA (Yang et al. 2018) dataset cites; confirm a Semantic
%   Chameleon / 2026 corpus-dependent-defense cite for Related Work positioning.
% =====================================================================================
\bibliographystyle{IEEEtran}
\begin{thebibliography}{99}
\bibitem{poisonedrag} W. Zou, R. Geng, B. Wang, and J. Jia, ``PoisonedRAG: Knowledge corruption attacks to retrieval-augmented generation of large language models,'' in \emph{Proc. 34th USENIX Security Symposium}, 2025.
\bibitem{avfilter} S. Choudhary, N. Palumbo, A. Hooda, K. Dvijotham, and S. Jha, ``Through the stealth lens: Rethinking attacks and defenses in RAG,'' 2025.
\bibitem{robustrag} C. Xiang, T. Wu, Z. Zhong, D. Wagner, D. Chen, and P. Mittal, ``Certifiably robust RAG against retrieval corruption,'' in \emph{Advances in Neural Information Processing Systems (NeurIPS)}, 2025.
\bibitem{ragdefender} M. Kim, H. Lee, and H. Koo, ``Rescuing the unpoisoned: Efficient defense against knowledge corruption attacks on RAG systems,'' in \emph{Annual Computer Security Applications Conference (ACSAC)}, 2025.
\bibitem{catpoison} H. Mei, Z. Du, C. Cai, C. Cheng, and J. Chen, ``CatPoison: Category-oriented knowledge poisoning attacks in retrieval-augmented generation systems,'' in \emph{IEEE TrustCom}, 2025.
\bibitem{autotpa} C. Jin, J. Ma, X. Chen, Y. Ma, and Z. Zhou, ``AutoTPA: Automated and efficient trigger-targeted data poisoning in retrieval-augmented generation,'' in \emph{IEEE ICPADS}, 2025.
\bibitem{ragshield} K. Patil, ``RAGShield: Provenance-verified defense-in-depth against knowledge base poisoning in government retrieval-augmented generation systems,'' 2026.
\bibitem{regain} S. Shajarian, K. Marsh, J. Benson, S. Khorsandroo, and M. Abdelsalam, ``ReGAIN: Retrieval-grounded AI framework for network traffic analysis,'' 2025.
\bibitem{bge} S. Xiao, Z. Liu, P. Zhang, and N. Muennighoff, ``C-Pack: Packaged resources to advance general Chinese embedding,'' 2023.
\bibitem{faiss} M. Douze, A. Guzhva, C. Deng, J. Johnson, G. Szilvasy, P.-E. Mazar\'{e}, M. Lomeli, L. Hosseini, and H. J\'{e}gou, ``The FAISS library,'' 2024.
\bibitem{hnsw} Y. A. Malkov and D. A. Yashunin, ``Efficient and robust approximate nearest neighbor search using hierarchical navigable small world graphs,'' \emph{IEEE Trans. Pattern Anal. Mach. Intell.}, vol. 42, no. 4, pp. 824--836, 2020.
\bibitem{wikitext} S. Merity, C. Xiong, J. Bradbury, and R. Socher, ``Pointer sentinel mixture models,'' in \emph{Int. Conf. on Learning Representations (ICLR)}, 2017.
\bibitem{gcg} A. Zou, Z. Wang, J. Z. Kolter, and M. Fredrikson, ``Universal and transferable adversarial attacks on aligned language models,'' 2023.
\bibitem{autodan} X. Liu, N. Xu, M. Chen, and C. Xiao, ``AutoDAN: Generating stealthy jailbreak prompts on aligned large language models,'' 2024.
\bibitem{shap} S. M. Lundberg and S.-I. Lee, ``A unified approach to interpreting model predictions,'' in \emph{Advances in Neural Information Processing Systems (NeurIPS)}, 2017, pp. 4768--4777.
\end{thebibliography}

\end{document}
